%% ma: L12-B08-0111
%% lop: 12
%% chuong: 2
%% bai: 8
%% dang: 12.8.25
%% loai: TN4
%% muc: TH
%% dap_an: "A"
%% nguon: "CĐ2.3.pdf giáo viên gửi 27/10/2026 (M07, Lê Minh Khôi – Nguyễn Anh Khôi, Chương II Chủ đề 3 Biểu thức tọa độ của các phép toán vectơ), Câu 23"
%% kien_thuc: "Bài 8 (hai vectơ cùng phương)"
%% trang_thai: chinh-thuc
\begin{ex}
Trong không gian $Oxyz$, cho hai điểm $M(3;1;-1)$ và $N(2;-1;4)$. Vectơ nào sau đây cùng phương với vectơ $\overrightarrow{MN}$?
\choice{\True $\vec a=(1;2;-5)$}{$\vec b=(4;8;-5)$}{$\vec c=(-1;-2;6)$}{$\vec d=(1;2;-6)$}
\loigiai{$\overrightarrow{MN}=(-1;-2;5)=-\vec a$ nên $\vec a$ cùng phương với $\overrightarrow{MN}$.}
\end{ex}
